We estimate the classical mass-weighted vibrational density of states from
the reference and MLIP trajectories using identical preprocessing and
spectral-estimation settings. Positions are first unwrapped across periodic
boundaries. Velocities are reconstructed by central differences,
$\mathbf{v}_i(t_j)=[\mathbf{r}_i(t_{j+1})-\mathbf{r}_i(t_{j-1})]/(2\Delta t)$,
discarding the two endpoint frames. We use Welch averaging of one-sided
Hann-windowed velocity periodograms, with 50\% overlap between segments.
The requested segment duration is 2~ps; for shorter matched trajectories,
it is reduced identically for the reference and MLIP data to obtain at
least four overlapping segments. The temporal mean of each Cartesian
velocity component is removed within each segment. Each atom's component
spectra are weighted by its ASE standard atomic mass, summed over atoms and Cartesian
directions, and averaged over segments. The resulting spectrum is
normalised to unit area before evaluating Eq.~(\ref{eq:vdos-error}).
The spectra displayed in the figures are the same spectra used for scoring.
Fourfold zero padding provides a denser frequency grid for integration and
plotting without increasing the physical spectral resolution. Actual
segment durations, segment counts and spectral bandwidths are recorded
with the results. We assess sensitivity to the requested segment duration,
minimum segment count, frequency grid and mass-weighting convention.
The short hydrogen and CuAu reference trajectories impose coarser spectral
resolution; spectral errors retain finite-trajectory sampling uncertainty.
Segment mean removal suppresses near-zero-frequency and diffusive power;
the resulting spectra quantify vibrational shape rather than diffusion.
Central differencing also attenuates high-frequency velocity power, with
the same frequency response applied to both trajectories.
